\clearpage\begingroup\NoSourceHead\SourceLeaf{457}
\setlength{\fboxsep}{5pt}
\noindent\fbox{\begin{minipage}{\dimexpr\textwidth-2\fboxsep-2\fboxrule\relax}
\noindent\fbox{\begin{minipage}{\dimexpr\linewidth-2\fboxsep-2\fboxrule\relax}
\fontsize{8}{9.5}\selectfont
\begin{center}
{\Large C. F. GAUSS WERKE}\\[5pt]
HERAUSGEGEBEN VON DER KGL GESELLSCHAFT DER WISSENSCHAFTEN\\
ZU GÖTTINGEN\\[3pt]
IN KOMMISSION BEI B G TEUBNER IN LEIPZIG\\[3pt]
\rule{.35em}{.2pt}\hspace{3.5em}\rule{.55em}{.2pt}\par\medskip
{\large A. AUSGABE AUF DRUCKPAPIER.}
\end{center}
\newcommand{\GC}[2]{\par\noindent\hangindent=3em\hangafter=1\makebox[3em][r]{#1\enspace}#2\par}
\noindent Band\quad I.\quad DISQUISITIONES ARITHMETICAE Zweiter Abdruck 1870
\GC{II}{HÖHERE ARITHMETIK. Zweiter Abdruck. 1876}
\GC{}{Nachträge zum ersten Abdruck des zweiten Bandes 1876}
\GC{III.}{ANALYSIS. Zweiter Abdruck 1876}
\GC{IV.}{WAHRSCHEINLICHKEITSRECHNUNG UND GEOMETRIE. Zweiter Abdruck 1880.}
\GC{V.}{MATHEMATISCHE PHYSIK. Zweiter Abdruck 1877}
\GC{VI.}{ASTRONOMISCHE ABHANDLUNGEN UND AUFSÄTZE 1874}
\GC{VII}{THEORIA MOTUS UND THEORETISCH-ASTRONOMISCHER NACHLASS. Parabolische Bewegung, Störungen der Ceres und der Pallas, Theorie des Mondes., 1906.}
\GC{VIII}{ARITHMETIK, ANALYSIS, WAHRSCHEINLICHKEITSRECHNUNG, GEOMETRIE (Nachträge zu Band I—IV.) 1900}
\GC{IX.}{GEODÄSIE. (Fortsetzung von Band IV ) 1903\raisebox{.65ex}{.}}
\GC{X}{Abteilung 1. NACHLASS UND BRIEFWECHSEL ZUR REINEN MATHEMATIK. (Nachträge zu Bd. I—IV und VIII. TAGEBUCH. 1917.}
\GC{XI.}{Abteilung 1. NACHLASS UND BRIEFWECHSEL ZUR PHYSIK, ASTRONOMIE UND CHRONOLOGIE. Unter der Presse.}
\GC{X.}{Abteilung 2. AUFSÄTZE ÜBER GAUSS’ WISSENSCHAFTLICHE TÄTIGKEIT AUF DEN GEBIETEN DER REINEN MATHEMATIK. Unter der Presse.}
\GC{XI.}{Abteilung 2. AUFSÄTZE ÜBER GAUSS’ WISSENSCHAFTLICHE TÄTIGKEIT AUF DEN GEBIETEN DER ANGEWANDTEN MATHEMATIK In Vorbereitung.}
\par\noindent\hspace*{3em}Die Bände X\textsubscript{2} und XI\textsubscript{2} bestehen aus einzelnen Aufsätzen ,Essays’, die besonders paginiert sind und in der Reihenfolge ihrer Fertigstellung auch gesondert ausgegeben werden. Mit dem letzten Aufsatz eines jeden Bandes wird Titelblatt und Inhaltsverzeichnis geliefert.
\GC{XII.}{BIOGRAPHISCHES NEBST EINEM GENERALREGISTER. In Vorbereitung.}
\begin{center}{\large B. AUSGABE AUF SCHREIBPAPIER.}\end{center}
\noindent Band I—VI (nur vollständig abzugeben). Band VII. Band VIII. Band IX. Band X. Abteilung 1.
\par\smallskip\begin{center}
\begin{tikzpicture}[x=.25pt,y=1pt,line width=.2pt]
\draw(0,0)--(9,0);\draw(14,0)--(24,0);\draw(34,0)--(44,0);\draw(69,0)--(77,0);
\draw(117,0)--(126,0);\draw(128,0)--(141,0);\draw(143,0)--(172,0);\draw(174,0)--(194,0);\draw(197,0)--(209,0);
\end{tikzpicture}\end{center}
\par\noindent\hspace*{\parindent}Berichte über den Stand der Herausgabe von \textsc{Gauss}’ Werken werden, von 1898 beginnend, regelmäßig in den Nachrichten der Kgl. Gesellschaft der Wissenschaften zu Göttingen (geschäftliche Mitteilungen) veröffentlicht.
\par\bigskip\noindent\rule{\linewidth}{.3pt}\par\vspace{-6pt}\noindent\rule{\linewidth}{.3pt}
\begin{center}\fontsize{6.5}{8}\selectfont GÖTTINGEN, GEDRUCKT IN DER DIETERICHSCHEN UNIV.-BUCHDRUCKEREI (W. FR. KAESTNER).\end{center}
\end{minipage}}\end{minipage}}
\clearpage\SourceLeaf{458}\NoSourceHead
\begin{center}
\vspace*{.7in}
{\Large GAUSS UND DIE VARIATIONSRECHNUNG}\\[.8in]
VON\\[.8in]
{\large OSKAR BOLZA}
\end{center}
\vfill\noindent{\small X\textsubscript{2} Abh. 5.}\hfill{\small 1}
\clearpage\endgroup
